\documentclass[../../main.tex]{subfiles}

\begin{document}

\section{自反性与弱紧性: Kakutani 定理}\label{section:ch15-kakutani}

\begin{proposition}\label{proposition:ch15-natural-homeo}
令 $X$ 为赋范线性空间. 则自然嵌入 $J:X\to X^{**}$ 是局部凸拓扑向量空间 $X$ 与 $J(X)$ 之间的拓扑同胚, 其中 $X$ 有弱拓扑而 $J(X)$ 有弱 $^{*}$ 拓扑.
\end{proposition}

\begin{proof}
令 $x_{0}$ 属于 $X$. 弱拓扑在 $x_{0}\in X$ 的邻域基定义为以下形式的集合: 对 $\varepsilon>0$ 与 $\psi_{1},\cdots,\psi_{n}\in X^{*}$,
\begin{align*}
\mathcal{N}_{\varepsilon,\psi_{1},\cdots,\psi_{n}}(x_{0})=\left\{x\in X\ \middle|\ \left|\psi_{i}(x-x_{0})\right|<\varepsilon,\ 1\leqslant i\leqslant n\right\}.
\end{align*}
现在对每个 $x\in X$ 与 $1\leqslant i\leqslant n$, $[J(x)-J(x_{0})]\psi_{i}=\psi_{i}(x-x_{0})$, 因此,
\begin{align*}
J(\mathcal{N}_{\varepsilon,\psi_{1},\cdots,\psi_{n}}(x_{0}))=J(X)\cap\mathcal{N}_{\varepsilon,\psi_{1},\cdots,\psi_{n}}(J(x_{0})).
\end{align*}
因此 $J$ 将 $X$ 中的弱拓扑在原点的基映上 $J(X)$ 中的弱 $^{*}$ 拓扑在原点的基. 因此 $J$ 是从具有弱拓扑的 $X$ 映上具有弱 $^{*}$ 拓扑的 $J(X)$ 的同胚.
\end{proof}

\begin{theorem}[Kakutani]\label{theorem:ch15-kakutani}
Banach 空间是自反的当且仅当它的闭单位球是弱紧的.
\end{theorem}

\begin{proof}
令 $X$ 为 Banach 空间. 记 $X$ 与 $X^{**}$ 的闭单位球分别为 $B$ 与 $B^{**}$. 假设 $X$ 是自反的. 自然嵌入是同构, 因此 $J$ 是将 $B$ 映上 $B^{**}$ 的一一映射. 另一方面, 根据命题\cref{proposition:ch15-natural-homeo}, $J$ 是从具有弱拓扑的 $B$ 到具有弱 $^{*}$ 拓扑的 $B^{**}$ 的同胚, 根据 Alaoglu 定理, 其中 $X$ 被 $X^{*}$ 代替, $B^{**}$ 是弱 $^{*}$ 紧的, 因此任何同胚于它的拓扑空间也是紧的. 特别地, $B$ 是弱紧的.

现在假设 $B$ 是弱紧的. 紧拓扑空间的连续像是紧的. 我们从命题\cref{proposition:ch15-natural-homeo}推出 $J(B)$ 关于弱 $^{*}$ 拓扑是紧的. 由于弱 $^{*}$ 拓扑是 Hausdorff 的, 因此 $J(B)$ 是闭的. 显然, $J(B)$ 是凸的. 为证明 $X$ 的自反性, 我们用反证法. 假设 $X$ 不是自反的. 令 $T$ 属于 $B^{**}\sim J(B)$. 应用 $X$ 被 $X^{*}$ 代替而 $W=J(X^{*})$ 时的超平面分离定理的推论 26, 因此存在泛函 $\psi\in X^{*}$, 使得 $\|\psi\|=1$ 且
\begin{align*}
T(\psi)<\inf_{S\in J(B)}S(\psi)=\inf_{x\in B}\psi(x).
\end{align*}
右边的下确界等于 $-1$, 由于 $\|\psi\|=1$. 因此 $T(\psi)<-1$. 这是一个矛盾, 由于 $\|T\|\leqslant 1$ 且 $\|\psi\|=1$. 因此 $X$ 是自反的.
\end{proof}

\begin{corollary}\label{corollary:ch15-closed-bounded-convex}
自反 Banach 空间的每个闭有界凸子集是弱紧的.
\end{corollary}

\begin{proof}
令 $X$ 为 Banach 空间. 根据 Kakutani 定理, $X$ 的闭单位球是弱紧的. 因此任何闭球也如此. 根据 Mazur 定理, $X$ 的每个闭凸子集是弱闭的. 因此 $X$ 的任何闭凸有界子集是弱紧集的弱闭子集, 因此必须是弱紧的.
\end{proof}

\begin{corollary}\label{corollary:ch15-reflexive-dual}
令 $X$ 为自反 Banach 空间. 则它的对偶空间的闭单位球 $B^{*}$ 关于弱 $^{*}$ 拓扑是序列紧的.
\end{corollary}

\begin{proof}
由于 $X$ 是自反的, $B^{*}$ 上的弱拓扑与弱 $^{*}$ 拓扑是相同的. 因此, 根据 Alaoglu 定理, $B^{*}$ 是弱紧的. 我们从 Kakutani 定理推出 $X^{*}$ 是自反的. 因此从前一章的定理 17 推出 $X^{*}$ 中的每个有界序列有弱 $^{*}$ 收敛子序列. 但 $B^{*}$ 是弱 $^{*}$ 闭的, 因此 $B^{*}$ 关于弱 $^{*}$ 拓扑是序列紧的.
\end{proof}

\begin{exercise}\label{exercise:ch15-06}
证明赋范线性空间的每个弱紧子集关于范数是有界的.
\end{exercise}

\begin{exercise}\label{exercise:ch15-07}
证明 Banach 空间 $X$ 的对偶 $X^{*}$ 的闭单位球 $B^{*}$ 有一个极值点.
\end{exercise}

\begin{exercise}\label{exercise:ch15-08}
令 $\mathcal{T}_{1}$ 与 $\mathcal{T}_{2}$ 为集合 $\mathcal{S}$ 上的两个紧 Hausdorff 拓扑, 满足 $\mathcal{T}_{1}\subseteq\mathcal{T}_{2}$. 证明 $\mathcal{T}_{1}=\mathcal{T}_{2}$.
\end{exercise}

\begin{exercise}\label{exercise:ch15-09}
令 $X$ 为包含子空间 $Y$ 的赋范线性空间. 对 $A\subseteq Y$, 证明 $A$ 上由 $Y^{*}$ 诱导的弱拓扑与 $A$ 承袭自 $X$ 的弱拓扑是相同的.
\end{exercise}

\begin{exercise}\label{exercise:ch15-10}
用以下步骤证明 Banach 空间 $X$ 是自反的当且仅当它的对偶空间 $X^{*}$ 是自反的.
\begin{enumerate}[(i)]
\item 若 $X$ 是自反的, 证明 $B^{*}$ 上的弱与弱 $^{*}$ 拓扑是相同的, 从这推出 $X^{*}$ 是自反的.
\item 若 $X^{*}$ 是自反的, 用(i)部分与第 14 章的命题 15 证明 $X$ 是自反的.
\end{enumerate}
\end{exercise}

\begin{exercise}\label{exercise:ch15-11}
对 Banach 空间 $X$, 根据前一个习题, 若 $X$ 是自反的, 则 $X^{*}$ 也是. 若存在 $X^{*}$ 的不是自反的闭子空间, 则 $X$ 不是自反的. 令 $K$ 为无穷紧 Hausdorff 空间而 $\{x_{n}\}$ 是 $K$ 的可数无穷子集的列举. 定义算子 $T:\ell^{1}\to[C(K)]^{*}$ 为
\begin{align*}
\text{对所有 }\{\eta_{k}\}\in\ell^{1}\text{ 与 }f\in C(K),\quad[T(\{\eta_{k}\})](f)=\sum_{k=1}^{\infty}\eta_{k}\cdot f(x_{k}).
\end{align*}
证明 $T$ 是等距. 因此, 由于 $\ell^{1}$ 不是自反的, $T(\ell^{1})$ 与 $C(K)$ 也不是. 用维数计算方法证明: 若 $K$ 是有限集, 则 $C(K)$ 是自反的.
\end{exercise}

\begin{exercise}\label{exercise:ch15-12}
若 $Y$ 是 Banach 空间 $X$ 的线性子空间, 我们定义零化子 $Y^{\perp}$ 为 $X^{*}$ 的使得在 $Y$ 上 $\psi=0$ 的那些 $\psi\in X^{*}$ 组成的子空间. 若 $Y$ 是 $X^{*}$ 的子空间, 我们定义 $Y^{0}$ 为 $X$ 中使得对所有 $\psi\in Y$ 满足 $\psi(x)=0$ 的向量的子空间.
\begin{enumerate}[(i)]
\item 证明 $Y^{\perp}$ 是 $X^{*}$ 的闭线性子空间.
\item 证明 $(Y^{\perp})^{0}=\overline{Y}$.
\item 若 $X$ 是自反的而 $Y$ 是 $X^{*}$ 的子空间, 证明 $Y^{\perp}=J(Y^{0})$.
\end{enumerate}
\end{exercise}

\end{document}